% !TeX root = ../thuthesis-example.tex

\begin{resolution}
论文围绕车云协同液罐车防侧翻规划与控制开展研究，选题具有重要理论意义和实用价值。

论文工作取得的成果如下：

（1）构建了车云协同液罐车防侧翻场景任务体系架构，解决了车路云复杂体系安全机制不明确的问题，为算法设计与工程部署提供了可追溯的约束依据；

（2）提出了面向交通交互场景的图扩散云端预测性决策与规划方法，解决了大时空尺度下车辆不确定性交互表征困难的问题，实现了长时域交互风险演化建模与诱发性风险规避；

（3）设计了罐液耦合高精度计算模型并设计模型预测控制器，解决了现有控制方法忽略横纵向耦合晃动导致模型在纵向激励下失配的问题，形成了长短时域协同的防侧翻保障。

论文层次清晰，文字通顺，图表规范。

论文工作表明戚笑景同学在本门学科及相关学科领域掌握了坚实的基础理论和系统的专门知识，具有从事科学研究和独立担负专门技术工作的能力。

答辩过程中，报告简明扼要，回答问题正确。经无记名投票，答辩委员会5票一致同意通过论文答辩，并建议授予戚笑景工学硕士学位。

\end{resolution}
